% ============================================================
% CHAPTER 5: DETECTION METHODOLOGY
%
% Structure (matches ToC):
%   5.1  Shared Measurement Conventions
%   5.2  Self-Reinforcement
%   5.3  Multistability
%   5.4  Dominance
%   5.5  Limit Cycles          (subsubsections collapsed)
%
% Every detector subsection follows the same five moves:
%   Concept -> Measurement -> Detection -> Worked example -> Limitations
% Split rule: Measurement is what we compute from the log; Detection is
% the null it is tested against and the rule that flags the motif.
%
% Notation is unified with Chapter 4:
%   t = round index, r = repetition, N = agents (=4), M = options
%   v_i^t = vote, c_i^t = confidence, s_t = macrostate.
% Recurrence-matrix positions are indexed by (a,b) to avoid clashing
% with the agent index i.
%
% REMOVED LABELS (redirect any \autoref pointing at them to
% \autoref{sec:method:lc}):
%   sec:method:lc:agent
%   sec:method:lc:system
%
% Persuasiveness and message alignment are retained (commented out at
% the end of this file) for later reference but are no longer part of
% the methodology. Per-detector implementation notes and the list of
% future motifs are moved to the appendix.
% ============================================================


Each of the four motifs defined in this section is a hypothesis about the
underlying debate dynamics. If the system self-organizes in the way a
complex-systems reading predicts, then confidence should reinforce within a held
position, repeated repetitions should settle into a small number of
initial-condition-selected attractors, influence should concentrate on particular
agents, and in some configurations vote sequences may oscillate persistently
without converging. This section turns each hypothesis into a \emph{detector}: a
statistic computed from the interaction log together with a principled null model
against which it is tested, so that a positive result reflects genuine structure
rather than the chance regularities of a small, discrete state space. Together the
four detectors form \suitename{}, the detector suite this thesis contributes.

Two properties are shared by all four detectors and therefore are worth stating up
front. Firstly, each detector reads only from ordinary interaction logs, the vote
and confidence channels defined in \autoref{sec:exp}, which is what later lets us
propose the same detectors as diagnostic instruments. And secondly, each detects a
\emph{signature}, not a certified dynamical object. Thus, we add a
``Limitations'' paragraph to each detector which makes explicit what the
corresponding statistic can and cannot claim.


% ------------------------------------------------------------
% 5.1 SHARED MEASUREMENT CONVENTIONS
% ------------------------------------------------------------
\subsection{Shared Measurement Conventions}
\label{sec:method:vocab}

Before defining each detector individually, it is worth fixing what they share.
Four things recur across every subsection below: the shape those subsections take,
the classification of the motifs themselves, the macrostate the system-level
detectors read from, and the nesting of the data that constrains every comparison.

The four subsections each make the same five moves, in the same order.
\emph{Concept} states the motif in dynamical-systems terms and names the quantity
in a debate that plays the role of the state variable. \emph{Measurement} defines
the \emph{measurement primitive}, the lowest-level object the detector scores, and
the statistic computed from it. \emph{Detection} gives the null model that
statistic is tested against and the rule that flags the motif. \emph{Worked
example} runs the detector by hand on a small case, and \emph{Limitations} states
what the resulting statistic cannot claim. The division between the second and
third move is the one worth holding on to, since it recurs throughout: measurement
is what we compute from the log, detection is what we compare it against. So that
the four worked examples can be read against one another, they share a single
running scenario, one task with $N = 4$ agents and $M = 4$ options labelled $A$ to
$D$ of which $A$ is correct, and they write an agent's state as its option with its
confidence as a subscript, so that $D_8$ is a vote for $D$ held at confidence $8$.

The motifs those subsections detect are preselected rather than exhaustively
enumerated, each motivated by a documented failure mode of LLM-based debate.
\autoref{tab:method:motifs} lists them together with their classification and the
section in which each detector is defined. Further motifs that fall within the same
structural-dynamical framework but lie outside the scope of this study are catalogued
in Appendix~\ref{app:future_motifs}.

\begin{table}[H]
\centering
\small
\begin{tabular}{llll}
\toprule
Motif & Classification & Level & Section \\
\midrule
Self-reinforcement & Structural (positive autoregulation) & Agent & \autoref{sec:method:sr} \\
Multistability        & Dynamical (coexisting attractors)    & System & \autoref{sec:method:multistability} \\
Dominance          & Structural (emergent hub)            & System & \autoref{sec:method:dominance} \\
Limit cycle        & Dynamical (periodic orbit)           & Agent \& System & \autoref{sec:method:lc} \\
\bottomrule
\end{tabular}
\caption[Motif classification]{The four motifs classified by type (structural or dynamical, following
\autoref{sec:bg:motifs}) and level (whether the detector reads from individual
agent trajectories, the collective macrostate, or both).}
\label{tab:method:motifs}
\end{table}

What the table calls the \emph{level} of a motif decides where its detector reads
from. Self-reinforcement, the only agent-level motif, reads directly from
per-agent vote and confidence trajectories. The three system-level readings
instead operate on the agent-anonymised macrostate $s_t$ defined in
\autoref{sec:setup:notation}, whose key property is that composition retains more
information than a plain vote count: a $3$--$1$ split and a $2$--$2$ split are
distinct macrostates, and options remain labelled because one is correct and they
are not interchangeable. Each of the three takes a different slice of that record,
however. Multistability reads only the locked-in endpoint, dominance reads the
emergent influence graph, which is built relationally and therefore needs
within-repetition agent identity, and the system-level limit-cycle test reads the
full sequence $s_0, \ldots, s_T$.

Whichever slice a detector takes, it inherits the same nesting of the data,
\begin{equation}
  \text{agent} \subset \text{repetition} \subset \text{task} \subset
  \text{config } (W \times \text{topology}) \subset \text{dataset},
\end{equation}
with $W \in \{1,2,5\}$, $\text{topology} \in \{\text{fc}, \text{star}\}$, $N = 4$,
and $\text{dataset} \in \{\text{GPQA}, \text{HiddenBench}\}$, and that nesting
constrains what may be compared with what. Because agents interact within a
repetition, their trajectories are not independent, so the repetition is the
independent unit of replication and agent-level trajectories are never pooled as
if they were separate observations (pseudoreplication). For the same reason GPQA
and HiddenBench are always analysed separately and never pooled, and memory-window
comparisons are made within a fixed topology and dataset, since cells from
different configs are not exchangeable.

A last property is shared by all four detectors, and it bounds every claim they
make: each is phenomenological rather than causal. We detect a trajectory or
influence \emph{signature}, not a certified attractor, cycle, or causal hub, and
separating the signature from its underlying cause requires dedicated perturbation
experiments, which we leave to future work. Under the cooperative,
convergence-forced design we also expect several detectors to return sparse or
near-null results, which is why each detection paragraph below names the baseline
its statistic is compared against rather than testing it against zero.


% ------------------------------------------------------------
% 5.2 SELF-REINFORCEMENT
% ------------------------------------------------------------
\subsection{Self-Reinforcement}
\label{sec:method:sr}

\paragraph{Concept.}
In complex systems, self-reinforcement is \emph{positive feedback}: the current
level of a state variable raises the drive toward more of the same, so deviations
are amplified rather than damped \citep{arthur1989competing}. At the motif level
this is positive autoregulation, a node that augments its own production
\citep{milo2002network, alon2007network}. For us the state variable is an agent's
expressed \emph{confidence}. Self-reinforcement is then the tendency of an agent
to push its own confidence upward while it holds a position, driven by
accumulating commitment rather than by what the incoming information warrants. The
pattern is already documented: in multi-turn LLM debate, confidence tends to rise
across rounds in an anti-Bayesian fashion, even where no new information justifies
the rise \citep{prasad2025certain}.

For the detection of self-reinforcement a rising confidence trajectory is on its
own necessary but not sufficient, since rational updating and peer conformity also
push confidence up. We therefore measure two signatures that go beyond it.
\emph{Directional drift} is the tendency of confidence to increase rather than
decrease while a position is held. \emph{Strength scaling} is the prediction that
this drift grows with how strongly an agent is re-exposed to its own prior
commitments, which the memory window $W$ controls directly. A third signature,
\emph{endogeneity}, would establish that the rise comes from the agent's own
commitment rather than from new information, but it cannot be separated from peer
effects using trajectory data alone.

\paragraph{Measurement.}
The primitive is the \emph{vote-stable segment}: a maximal block of consecutive
rounds over which agent $i$'s vote does not change, running from its start $t_s$
to its end $t_e$, so that $v_i^t = v_i^{t_s}$ for all $t \in \{t_s, \ldots, t_e\}$.
Writing $L = t_e - t_s + 1$ for its length in rounds, a segment is retained only if
$L \geq 3$. Two rounds would already define a slope, but that slope would be a
single confidence difference with no residual degrees of freedom, so its sign would
turn on one integer step; three rounds give two steps for the least-squares fit to
average over. An agent might contribute several segments if it flips and one long
segment if it locks in early. Agents are exchangeable, so segments are pooled
within a config and never matched across repetitions.
Within a retained segment we regress confidence on the round index by ordinary
least squares,
\begin{equation}
  c_i^t = \beta_0 + \beta_{\mathrm{seg}}\, t + \varepsilon_t,
  \qquad t \in \{t_s, \ldots, t_e\},
\end{equation}
and keep the slope $\beta_{\mathrm{seg}}$, in confidence points per round. Its sign
feeds the prevalence test below, and its signed value feeds the mean slope
$\bar{\beta}$ used in the $W$ comparison.

\paragraph{Detection.}
With the segments in hand, the question is whether their slopes lean upward more
often than chance allows. We read them in two ways, both conditioned on the
config.
The first is a presence test for directional drift. Writing $n_{\mathrm{seg}}$ for
the number of retained segments in a config-task pair, the prevalence of upward
segments is
\begin{equation}
  \hat{p}_{\mathrm{SR}}
  = \frac{\bigl|\{\text{seg}: \beta_{\mathrm{seg}} > 0\}\bigr|}{n_{\mathrm{seg}}},
\end{equation}
so flat segments ($\beta_{\mathrm{seg}} = 0$) sit in the denominator but never in
the numerator. Under a no-feedback ``wobble'' null an upward and a downward
segment are equally likely, so with a fraction $f$ of flat segments
$P(\beta_{\mathrm{seg}} > 0) = \tfrac{1}{2}(1-f) \leq 0.5$. Thus, self-reinforcement is
indicated when $\hat{p}_{\mathrm{SR}} > 0.5$. Both the flat segments and the
confidence ceiling, since a segment already at the top of the scale cannot drift
further up, make a positive result conservative.
The second reading uses the memory window to discriminate between mechanisms.
Letting $\bar{\beta}$ be the mean segment slope in a config, an \emph{endogenous}
reading predicts that re-reading one's own prior commitments across a longer
window amplifies the internal drive,
\begin{equation}
  \bar{\beta}(W{=}5) > \bar{\beta}(W{=}2) > \bar{\beta}(W{=}1),
\end{equation}
whereas under a \emph{peer-conformity} reading a longer window mainly adds peer
content, and the ordering may be absent or reversed. 

Both readings are reported per config-task and then summarised across tasks,
for example ``above $0.5$ in $k$ of $50$ tasks'', so that an effect is shown to be
broad rather than driven by one task.

\paragraph{Worked example.}
In the running scenario, suppose agent~$0$ holds option~$A$ for five rounds with
confidences $c = (5, 6, 6, 8, 9)$. This is one vote-stable segment of length
$L = 5$, and regressing confidence on the round index gives
$\beta_{\mathrm{seg}} = +1.0$ confidence points per round: an upward segment,
counting toward $\hat{p}_{\mathrm{SR}}$. Had the agent instead drifted
$c = (7, 7, 6, 6, 5)$, the slope would be $-0.5$ and the segment would count
against it. A segment that never moved, $c = (6,6,6,6,6)$, has
$\beta_{\mathrm{seg}} = 0$: retained in the denominator but not counted as upward,
which is what makes the $0.5$ reference conservative. Under the strength-scaling
reading, the average of such slopes should grow as the window widens from
$W{=}1$ to $W{=}5$.

\paragraph{Limitations.}
Beyond the shared phenomenological caveat, two points bound the claim.
\emph{Endogeneity is not isolated}: an upward segment is equally consistent with
peer copying and with a peer surfacing genuinely new information in a later round,
and rising confidence alone cannot tell these apart. The $W$ ordering narrows the
ambiguity but does not settle it, since a longer window also re-exposes more peer
content. \emph{Convergence is a confound}: on easier tasks fast agreement can
itself raise confidence, which is why prevalence is reported per task and never as
a single pooled figure.


% ------------------------------------------------------------
% 5.3 MULTISTABILITY
% ------------------------------------------------------------
\subsection{Multistability}
\label{sec:method:multistability}

\paragraph{Concept.}
In dynamical systems, \emph{multistability} is the coexistence of two or more
stable attractors for one fixed parameter set. Into which attractor the system settles
into is decided by the basin of attraction containing the initial condition
\citep{feudel2008complex, pisarchik2014control}. At the motif level this is the dynamical signature of positive
feedback with a symmetry: the same interaction rules admit several self-consistent
outcomes \citep{milo2002network, alon2007network}. For us the state variable is
the macrostate $s_t$, and multistability is read off the \emph{locked-in endpoint}:
for a fixed task, the debate converges to different consensus answers on different
repetitions, and which one is reached is selected by the initial vote
configuration rather than by the task's correct option. This is the structural
reading of unreliable correction: on a task where the model carries a systematic
bias, a wrong-consensus and a right-consensus attractor can coexist, and the
group's fate is basin competition decided at round $0$, which self-reinforcement
then amplifies and locks in. For the motif detection, seeing more than one final answer is necessary but
not sufficient, as a task may simply be noisy. Thus, genuine multistability
additionally requires that the initial condition \emph{predicts} the endpoint
\citep{pisarchik2014control}.

\paragraph{Measurement.}
The primitive and independent unit is the repetition, characterised by its initial
vote composition $s_0 = (n_1^0,\ldots,n_M^0)$ and its final macrostate. A
repetition is converged if its final macrostate is unanimous, regardless
of whether that unanimity was reached via early stopping or at the round cap $T$.
We restrict the multistability analysis to converged repetitions because a stable
non-unanimous endpoint reflects a failure to reach collective agreement rather than
a genuine attractor. The endpoint attractor $A^{(r)}$ of a converged repetition is
the winning option, so there are at most $M$ distinct attractors per task.
Repetitions that end in a non-unanimous state are excluded.
Scoring is per task-config cell with $n_c$ denoting the number of converged
repetitions in that cell. The cell's \emph{dominant attractor} is the most
frequent endpoint option, and each repetition's start is summarised by
\begin{equation}
  g(s_0) \in \{0, 1, \ldots, N\},
\end{equation}
the number of agents whose round-$0$ vote is on the dominant option.
From these we compute the empirical basin distribution over endpoint attractors,
$\hat{P}(A = k) = \tfrac{1}{n_c}\sum_r \mathbf{1}[A^{(r)} = k]$, a basin-stability
estimate of relative basin volumes under the model's own initial distribution
\citep{menck2013basin}. The effective number of coexisting attractors is the
inverse Simpson (Hill) index \citep{hill1973diversity},
\begin{equation}
  N_{\mathrm{eff}} = \frac{1}{\sum_k \hat{P}(A = k)^2} \in [1, M].
\end{equation}
$N_{\mathrm{eff}} \approx 1$ is a single attractor, $\approx 2$ two coexisting
attractors, and so on, so $N_{\mathrm{eff}}$ reports \emph{how many} attractors
coexist. Because attractors are unanimous options, it is bounded by $M$. Since $M$
varies across tasks, we report $N_{\mathrm{eff}}$ raw and compare it only within a
fixed $M$. In practice GPQA is entirely $M = 4$; only HiddenBench
mixes $M = 3$ and $M = 4$, so its $N_{\mathrm{eff}}$ summaries are reported
separately by option count.

\paragraph{Detection.}
Coexistence alone does not establish multistability; the initial condition must
also \emph{select} the attractor \citep{feudel2008complex, pisarchik2014control},
and it is that association which carries the inferential claim. We measure it
between $g(s_0)$ and the endpoint $A$ with Cram\'er's $V$
\citep{cramer1946mathematical} on their contingency table,
\begin{equation}
  V = \sqrt{\frac{\chi^2 / n_c}{\min(R_{\mathrm{tab}} - 1,\, C_{\mathrm{tab}} - 1)}} \in [0, 1],
\end{equation}
where $\chi^2$ is the Pearson statistic and $R_{\mathrm{tab}}, C_{\mathrm{tab}}$
are the table's numbers of rows (distinct $g(s_0)$, at most $N+1$) and columns
(distinct endpoints, at most $M$). $V \approx 0$ means the start is irrelevant;
$V$ near $1$ means the start deterministically selects the endpoint. The null
shuffles the pairing of $g(s_0)$ and $A$ across the $n_c$ repetitions, preserving
both marginals while destroying the start-to-endpoint link
\citep{theiler1992surrogate}; for $B = 1000$ surrogates the add-one one-sided
$p$-value flags basin selection when $p < 0.05$.

The two readings combine into a single cell label,
\begin{equation}
  \text{label}(c) =
  \begin{cases}
    \text{monostable} & N_{\mathrm{eff}} < 1.5,\\[2pt]
    \text{multistable} & N_{\mathrm{eff}} \geq 1.5 \ \wedge\ p < 0.05,\\[2pt]
    \text{stochastic (not multistable)} & N_{\mathrm{eff}} \geq 1.5 \ \wedge\ p \geq 0.05,
  \end{cases}
\end{equation}
with the degree of multistability read from $N_{\mathrm{eff}}$ and capped at $M$.
The cutoff $1.5$ is a reporting threshold, not a test; the inferential claim rests
on $p$. The third case is the essential guard: coexistence without basin selection
is scatter, not multistability. Because per-cell power is limited, results lean on
the across-task distribution of $N_{\mathrm{eff}}$ and $V$ rather than any single
cell.

\paragraph{Worked example.}
In the running scenario, take a cell with $n_c = 30$ converged repetitions whose
dominant attractor is option~$A$. The coexistence reading depends only on the
endpoint distribution, and the label depends on whether the start predicts it.
\begin{itemize}
  \item \textbf{Monostable.} Endpoints $28 \times A$, $2 \times B$ give
        $\hat{P} = (0.933, 0.067)$ and
        $N_{\mathrm{eff}} = 1/(0.933^2 + 0.067^2) \approx 1.14 < 1.5$: one valley.
  \item \textbf{Multistable.} Endpoints $15 \times A$, $15 \times B$ give
        $N_{\mathrm{eff}} = 2.0$. If repetitions starting with more agents on $A$
        (high $g(s_0)$) end on $A$ and those starting with few end on $B$, then
        $V \approx 1$, $p < 0.05$: two basins selected by the start.
  \item \textbf{Tristable.} Endpoints $10 \times A$, $10 \times B$, $10 \times C$
        with the start predicting each give $N_{\mathrm{eff}} \approx 3.0$,
        $p < 0.05$ (the maximal degree for $M = 3$).
  \item \textbf{Stochastic.} Endpoints $15 \times A$, $15 \times B$ again give
        $N_{\mathrm{eff}} = 2.0$, but if $g(s_0)$ is unrelated to the endpoint then
        $V \approx 0$, $p \gg 0.05$: coexistence without basin selection, correctly
        \emph{not} labelled multistable.
\end{itemize}

\paragraph{Limitations.}
First,
classical basin stability draws initial conditions uniformly
\citep{menck2013basin}, whereas our starts come from the model's own round-$0$
distribution, so $\hat{P}(A)$ estimates \emph{effective} basin volume under that
distribution, not geometric volume. Second, because $g(s_0)$ counts agents already
on the eventual dominant option, a high $V$ is consistent both with genuine basin
selection and with ordinary majority-vote persistence; the test detects
start-to-endpoint dependence but does not by itself separate these two mechanisms,
so we read the association as necessary but not mechanistically decisive.


% ------------------------------------------------------------
% 5.4 DOMINANCE
% ------------------------------------------------------------
\subsection{Dominance}
\label{sec:method:dominance}

\paragraph{Concept.}
Dominance is the concentration of directed influence onto a single node, an
\emph{emergent hub} which is a structural motif \citep{milo2002network, alon2007network}
read at the system level, since it is a property of the whole influence
distribution. Its formal home is opinion dynamics. In the DeGroot model
\citep{degroot1974consensus} and its social-power extension, an agent's influence
equals its network centrality, and the group is \emph{autocratic} when these are
maximally nonuniform and \emph{democratic} when uniform \citep{jia2015opinion}.
Dominance is the autocratic side. For us the state variable is the distribution of
influence across the four agents, and dominance is the tendency of that
distribution to concentrate on one of them rather than to spread evenly. When it
does, a single overconfident agent can pull the rest of the group toward its answer
\citep{chen2023consensus, cho2025herd}, degrading collective accuracy even when
the majority would otherwise be correct \citep{wynn2025talk}. To avoid attributing concentration that is simply an artefact of the wiring, we
measure influence for each topology separately. The result is a
concentration signature rather than proof of causation, since an adoption edge
cannot distinguish genuine influence from coincidental agreement.

\paragraph{Measurement.}
The primitive and independent unit is the per-repetition influence graph, which
requires within-repetition agent identity.
Dominance is defined over agents and their conversions, so it is independent of
the option count $M$. A \emph{conversion} occurs when agent $j$ adopts, at
round $t{+}1$, a vote that a disagreeing neighbour $i$ held at round $t$. The
directed edge $i \to j$ accumulates the confidence-weighted conversions of $j$ onto
$i$'s vote, with credit split among the neighbours already holding the target vote,
\begin{equation}
  w_{ij} = \sum_{t} \frac{1}{m_j^t}\,
           \mathbf{1}\!\left[v_j^{t+1} = v_i^{t} \neq v_j^{t}\right] c_j^{t},
  \qquad
  m_j^t = \bigl|\{k \in \mathcal{N}(j) : v_k^{t} = v_j^{t+1}\}\bigr|,
\end{equation}
where $c_j^t$ is the converting agent's confidence, $\mathcal{N}(j)$ its
topological neighbourhood, and the factor $1/m_j^t$ shares credit equally among
all eligible sources. Per-agent influence is the weighted out-strength
$\sigma_i = \sum_j w_{ij}$. We use out-strength rather than eigenvector centrality
\citep{bonacich1972factoring}, which is undefined or degenerate on
non-strongly-connected and acyclic graphs \citep{newman2010networks} and
hub-localises \citep{martin2014localization}. Our graphs are tiny and sparse, 
making out-strength the only always-defined node score.

The statistic is the concentration of the resulting influence-share vector
$p_i = \sigma_i / \sum_k \sigma_k$, measured by the normalised Herfindahl index
\citep{hirschman1964paternity},
\begin{equation}
  D = \frac{\sum_i p_i^2 - \tfrac{1}{N}}{1 - \tfrac{1}{N}} \in [0, 1].
\end{equation}
$D = 0$ is democratic as the agents have equal influence, and $D = 1$ is a single pure hub. $D$ is
undefined when $\sum_k \sigma_k = 0$, that is when no conversion occurs at all,
and such repetitions are excluded.

\paragraph{Detection.}
A high $D$ in the star is uninformative on its own, since the centre is the only
agent a leaf can be converted by and the wiring alone forces concentration. The
null is therefore keyed to topology: it asks how concentrated $D$ would be if, on
this same wiring and these same conversion events, influence were assigned at
random. We use a constrained shuffle in the spirit of Theiler's ``Algorithm~0''
\citep{theiler1992surrogate} that keeps the converting agent, the round, the
weight, and the number of eligible sources fixed, and redraws only the source
\emph{identity}, uniformly among $\mathcal{N}(j)$.
In the fully-connected topology every agent has three neighbours, so the shuffle can reassign any
conversion to any of them and the resulting null is symmetric; concentration above
it is emergent. However, for the star topology a leaf has exactly one neighbour, so a leaf
conversion has only one possible source and the shuffle leaves it unchanged. Where
a repetition is made up of such conversions, every surrogate reproduces the
observed graph and the test can never flag. This makes the star null degenerate, so emergent concentration is identified only in
the fully-connected topology. Raw $D$ is still reported for both, but in the star
it describes the wiring rather than the behaviour.

A repetition is concentrated when its observed $D$ falls in the upper tail of its
own surrogate distribution. With $B = 1000$ surrogates the add-one one-sided
$p$-value is
\begin{equation}
  p = \frac{1 + \#\{\,b : D^{(b)} \geq D\,\}}{1 + B},
\end{equation}
and we flag the repetition when $p < 0.05$. For the fully-connected topology we
report the standardised excess
$z_r = (D_r - \mu_{\mathrm{null}})/\sigma_{\mathrm{null}}$ and the flag rate
against the $5\,\%$ the threshold produces by chance, each summarised per task and
then counted across tasks. As for self-reinforcement, an entrenching hub predicts
that concentration strengthens as the memory window widens at fixed topology, and
we separately report whether the hub's vote matches the final consensus and how
often it capitulates.

\paragraph{Worked example.}
In the running scenario with the fully-connected topology, agent~$0$ opens at
$D_8$ while agents~$1$, $2$, and $3$ open at $A_6$, $B_5$, and $C_5$
respectively. The vote sequence is
\[
\begin{array}{c|cccc}
t & \text{agent}_0 & \text{agent}_1 & \text{agent}_2 & \text{agent}_3\\\hline
0 & D_8 & A_6 & B_5 & C_5\\
1 & D_9 & D_6 & B_5 & C_5\\
2 & D_9 & D_7 & D_5 & C_5\\
3 & D_9 & D_7 & D_6 & D_5
\end{array}
\]
Three conversions occur. At round~$1$, agent~$1$ flips to $D$; only agent~$0$
held $D$ at round~$0$, so the full weight of agent~$1$'s confidence goes to
edge $0 \to 1$, giving $w_{01} = 6$. At round~$2$, agent~$2$ flips to $D$; both
agents~$0$ and $1$ held $D$ at round~$1$, so credit splits equally, giving
$w_{02} = w_{12} = 2.5$. At round~$3$, agent~$3$ flips to $D$; all three of
agents~$0$, $1$, and $2$ held $D$ at round~$2$, so credit splits three ways,
giving $w_{03} = w_{13} = w_{23} = \tfrac{5}{3}$.
The out-strengths are therefore $\sigma = (\tfrac{61}{6},\, \tfrac{25}{6},\,
\tfrac{5}{3},\, 0)$, and the influence shares $p \approx (0.635,\, 0.260,\, 0.104,\,
0)$, and $D = (0.482 - 0.25)/0.75 \approx 0.31$. Agent~$0$ accounts for $64\,\%$
of influence yet $D$ is well below $1$ because the credit-splitting rule makes
each early adopter a co-source for all later conversions, so staggered convergence
never produces a pure hub. $D = 1$ would require every conversion to trace back to
a single agent, which is only possible when every other agent flips simultanously.

\paragraph{Limitations.}
First, with four agents the four shares and a discrete null bound the minimum
$p$, so per-repetition power is low and the signal lives in the aggregate. Second,
a single-conversion repetition gives $D = 1$ trivially, which is controlled by
aggregation and the topology null rather than by discarding. Third, as noted
above, the degeneracy of the star null means the emergent-excess test is
identified only in the fully-connected topology, so any claim about emergent
concentration in the star rests on raw $D$ and its $W$ ordering rather than on the
null comparison.


% ------------------------------------------------------------
% 5.5 LIMIT CYCLES
% ------------------------------------------------------------
\subsection{Limit Cycles}
\label{sec:method:lc}

\paragraph{Concept.}
In dynamical systems a limit cycle is an isolated closed orbit, a self-sustained
periodic oscillation that requires no external forcing
\citep{strogatz2015nonlinear}. At the motif level, sustained oscillation is the
dynamical signature of feedback loops, which is the complementary case to the positive
autoregulation behind self-reinforcement \citep{milo2002network, alon2007network}.
This is the only one of our selected motifs that lives at both agent and system level, 
so there are two state variables
rather than one. At the agent level the state variable is the agent's vote
$v_i^t \in \{1,\ldots,M\}$, and
a limit cycle is an agent that keeps cycling through votes (for example
$A, B, A, B, \ldots$) without ever settling. At the system level it is the
macrostate $s_t$, and a limit cycle is a collective configuration that cycles
without settling. Neither implies the other, so both are detected and both are
reported. This matters because convergence to consensus is one system fate,
whereas a sustained oscillation is a qualitatively different one, periodic
disagreement that never resolves. Detecting it tests whether cooperative debate can
become dynamically trapped in a cycle rather than at a point. With only a handful
of symbols in our setting, states repeat by chance and a transient fluctuation that later converges is
not a cycle, so the definition requires both that the oscillation is
\emph{self-sustained}, so that the sequence does not converge, and that its periodic
structure \emph{beats a chance baseline}. We measure both, though the isolation and
stability of a true limit cycle cannot be certified from one observed trajectory
and are not claimed.

\paragraph{Measurement.}
The primitive depends on the level. At the agent level it is the per-agent vote
sequence within a repetition, $v = (v_0, v_1, \ldots, v_T)$, taking the committed
\updatephase{} vote at each round. At the system level it is the per-repetition
composition sequence $s = (s_0, s_1, \ldots, s_T)$, in which case the repetition is
both the measurement primitive and the independent unit. Everything that follows
applies unchanged to either, so we write $x = (x_0, \ldots, x_{L-1})$ for whichever
symbolic sequence is being scored.

A sequence that converges is not a limit cycle, so converged sequences are excluded
before anything is computed. A sequence counts as converged when it ends in a
constant run of at least $u = 3$ symbols, the settling criterion already used in
\autoref{sec:method:multistability}. The judgement is made on the sequence itself
rather than on the debate, which matters twice over: an agent can settle while the
debate around it continues, and a debate that deadlocks at a $3$--$1$ or $2$--$2$
split holds that composition to the cap $T$, which is a fixed point even though no
consensus was reached. Whatever survives is a limit-cycle candidate,
subject to a guard $L \geq 4$ that leaves room for one repeat of the shortest
possible, period-2, cycle. Converged sequences are dropped whole rather than
trimmed, since a trailing constant run cannot be separated into cycle-internal
repeats and genuine convergence.

Periodicity in the surviving candidates is quantified with categorical recurrence
quantification analysis \citep{cocodale2014crqa}. For sequence positions $a, b$ the
recurrence matrix records where the sequence revisits a symbol,
\begin{equation}
  R(a,b) = \mathbb{1}[\,x_a = x_b\,], \qquad a \neq b,
\end{equation}
excluding the main diagonal, which is trivially all ones
\citep{eckmann1987recurrence}. The pattern of ones distinguishes the cases:
a periodic orbit produces diagonal lines parallel to the main diagonal and equally
spaced by the period, a fixed point produces a solid block, and scattered points
mean no periodicity at all \citep{marwan2007recurrence}. Two quantities summarise
it. \emph{Determinism} is the fraction of recurrence points lying on diagonal lines
of length $l \geq l_{\min} = 2$ \citep{zbilut1992embeddings, webber1994dynamical},
\begin{equation}
  DET = \frac{\sum_{l \geq l_{\min}} l\,P(l)}{\sum_{l \geq 1} l\,P(l)},
\end{equation}
where $P(l)$ is the histogram of diagonal-line lengths. The \emph{dominant period}
is the lag at which recurrence is densest,
\begin{equation}
  RR(\tau) = \frac{1}{L-\tau}\sum_{a=0}^{L-1-\tau} R(a, a+\tau),
  \qquad \hat{P} = \arg\max_{\tau \geq 1} RR(\tau).
\end{equation}
Both are needed, because $DET$ alone cannot separate persistence from periodicity:
a solid block and equally-spaced diagonals inflate it equally, so a single stray
symbol in an otherwise constant sequence can reach a high $DET$ while oscillating
not at all. We therefore require $\hat{P} \geq 2$, and cap it at
$\hat{P} \leq \lfloor L/2 \rfloor$, since a longer period cannot recur twice
within the window.

\paragraph{Detection.}
The null uses random-shuffle surrogates, that is,
random permutations of the sequence's own symbol multiset
\citep{theiler1992surrogate}. Shuffling destroys temporal order while preserving
the symbol frequencies, so the chance-match rate implied by the alphabet is built
into the null automatically. For $B = 1000$ surrogates we compute $DET^{(b)}$ and
the add-one one-sided $p$-value \citep{phipson2010permutation}
\begin{equation}
  p = \frac{1 + \#\{\,b : DET^{(b)} \geq DET\,\}}{1 + B},
\end{equation}
and flag a limit cycle when
\begin{equation}
  p < 0.05 \quad\text{and}\quad 2 \leq \hat{P} \leq \lfloor L/2 \rfloor .
\end{equation}
We report a binary flag per candidate, prevalence per config-task as the flagged
share among eligible candidates, and summaries across tasks per topology and $W$.

The two levels share this machinery and differ only in power. The macrostate space
has $\binom{N + M - 1}{M - 1}$ compositions, $15$ for $M = 3$ and $35$ for
$M = 4$, against only $M$ raw vote symbols, so chance recurrence is lower and
$DET$ cleaner at the system level. Against that, there is one sequence per
repetition rather than one per agent, so candidates are far fewer.

\paragraph{Worked example.}
In the running scenario, suppose agent~$2$ fails to settle, cycling
$v = (A, B, A, B, A, B)$ across $L = 6$ rounds. It never ends in a constant run, so
it is a candidate, and $L \geq 4$. The recurrence matrix (main diagonal excluded)
is
\[
  R = \begin{pmatrix}
    \cdot & 0 & 1 & 0 & 1 & 0\\
    0 & \cdot & 0 & 1 & 0 & 1\\
    1 & 0 & \cdot & 0 & 1 & 0\\
    0 & 1 & 0 & \cdot & 0 & 1\\
    1 & 0 & 1 & 0 & \cdot & 0\\
    0 & 1 & 0 & 1 & 0 & \cdot
  \end{pmatrix}.
\]
The ones lie on diagonals equally spaced by $2$, giving $RR(1) = 0$, $RR(2) = 1$
and $RR(3) = 0$, so $\hat{P} = 2$, which satisfies the guard
$2 \leq \hat{P} \leq \lfloor 6/2 \rfloor = 3$. The offset-2 diagonal carries a line
of length $4$ and the offset-4 diagonal a line of length $2$, so
$DET = (4 + 2)/(4 + 2) = 1.0$.
The null shows why length matters. The multiset $\{3A, 3B\}$ has
$\binom{6}{3} = 20$ arrangements, and only the two perfectly alternating ones reach
$DET = 1$, so $p = 2/20 = 0.10$: at $L = 6$ even a perfect period-2 cycle cannot
clear $0.05$. The same pattern over a full-length candidate ($L \approx 16$) is
flagged decisively, so only long, non-converged sequences are informative. The two
exclusions are equally concrete: $(A,B,A,B,B,B,B)$ ends in a
constant run and is dropped as a fixed point, while $(A,B,B,B,B,B,B)$ reaches a
high $DET$ but has $\hat{P} = 1$ and is rejected by the period guard as
persistence rather than oscillation.
The same scenario shows why the two levels are reported separately. If agents~$0$
and $1$ swap $A \leftrightarrow B$ each round while agents~$2$ and $3$ hold their
votes, every $s_t$ is identical. The agent-level test flags two period-2 cycles
while the system-level test sees a fixed point and excludes the repetition
entirely. Conversely, a rotating majority can cycle $s_t$ with no single agent
being periodic.

\paragraph{Limitations.}
First, the alphabet is small at the agent level, which inflates
the chance recurrence rate, and at both levels the sequence length caps the
detectable period at $\lfloor L/2 \rfloor$. Second, the shuffle null controls the
chance-recurrence inflation, but the period guard makes the test conservative, so a
small nonzero flag rate remains under a chance baseline. We compare observed
prevalence against this floor rather than against zero.